% =====================================================================
%  preambulo.tex — EduCurator AI V1
%  Configuración común del documento. No contiene texto del proyecto.
% =====================================================================

% --- Idioma y codificación ------------------------------------------
\usepackage[T1]{fontenc}
\usepackage[utf8]{inputenc}
\usepackage[spanish,es-tabla,es-noquoting]{babel}
\usepackage{lmodern}
\usepackage{microtype}

% --- Geometría y espaciado ------------------------------------------
\usepackage[a4paper,top=2.8cm,bottom=2.8cm,left=3cm,right=2.5cm]{geometry}
\usepackage{setspace}
\onehalfspacing
\setlength{\parskip}{0.4em}

% --- Tablas, figuras y listas ---------------------------------------
\usepackage{booktabs}
\usepackage{longtable}
\usepackage{tabularx}
\usepackage{array}
\usepackage{multirow}
\usepackage{graphicx}
\usepackage{float}
\usepackage{enumitem}
\setlist{itemsep=0.2em,topsep=0.3em}

% Columna X justificada a la izquierda y columna de ancho fijo centrada
\newcolumntype{L}[1]{>{\raggedright\arraybackslash}p{#1}}
\newcolumntype{C}[1]{>{\centering\arraybackslash}p{#1}}

% --- Color y cajas ---------------------------------------------------
\usepackage{xcolor}
\definecolor{ecAzul}{HTML}{1F4E79}
\definecolor{ecGris}{HTML}{4D4D4D}
\definecolor{ecFondo}{HTML}{F2F5F8}
\definecolor{ecBorde}{HTML}{C3D3E0}

\usepackage[most]{tcolorbox}

% Caja para advertencias metodológicas (lo que NO se puede afirmar)
\newtcolorbox{advertencia}[1][]{%
  enhanced, breakable,
  colback=ecFondo, colframe=ecAzul,
  boxrule=0.4pt, left=6pt, right=6pt, top=5pt, bottom=5pt,
  fonttitle=\bfseries\small, coltitle=white,
  attach boxed title to top left={yshift=-2mm,xshift=4mm},
  boxed title style={colback=ecAzul, boxrule=0pt, arc=1pt},
  title={Advertencia metodológica}, #1}

% Caja para delimitar el estado de una afirmación
\newtcolorbox{estado}[1][]{%
  enhanced, breakable,
  colback=white, colframe=ecBorde,
  boxrule=0.6pt, left=6pt, right=6pt, top=5pt, bottom=5pt,
  fonttitle=\bfseries\small, coltitle=ecAzul,
  #1}

% --- Diagramas -------------------------------------------------------
\usepackage{tikz}
\usetikzlibrary{arrows.meta,positioning,shapes.geometric,fit,backgrounds,calc}

\tikzset{
  ecnodo/.style={draw=ecAzul, fill=ecFondo, rounded corners=2pt,
                 align=center, font=\small, inner sep=5pt, minimum height=8mm},
  ecvar/.style={draw=ecGris, fill=white, rounded corners=2pt,
                align=center, font=\small, inner sep=5pt, minimum height=8mm},
  ecflecha/.style={-{Stealth[length=2mm]}, draw=ecGris, thick}
}

% --- Encabezados y pies ---------------------------------------------
\usepackage{fancyhdr}
\pagestyle{fancy}
\fancyhf{}
\fancyhead[L]{\small\textcolor{ecGris}{EduCurator AI --- Investigación en curso}}
\fancyhead[R]{\small\textcolor{ecGris}{\thepage}}
\renewcommand{\headrulewidth}{0.3pt}

% --- Títulos ---------------------------------------------------------
\usepackage{titlesec}
\titleformat{\part}[display]
  {\normalfont\Large\bfseries\color{ecAzul}}{\partname\ \thepart}{12pt}{\Huge}
\titleformat{\section}
  {\normalfont\large\bfseries\color{ecAzul}}{\thesection}{1em}{}
\titleformat{\subsection}
  {\normalfont\normalsize\bfseries\color{ecGris}}{\thesubsection}{1em}{}

% --- Citas textuales -------------------------------------------------
\usepackage[autostyle=true,english=american]{csquotes}

% --- Bibliografía ----------------------------------------------------
% APA 7 con biblatex + biber. En Overleaf: seleccionar compilador
% pdfLaTeX y biber (Menu > Settings > TeX Live version reciente).
\usepackage[
  backend=biber,
  style=apa,
  sortcites=true,
  sorting=nyt,
  backref=false,
  uniquename=false
]{biblatex}
\DeclareLanguageMapping{spanish}{spanish-apa}
\addbibresource{referencias.bib}

% --- Hipervínculos (cargar al final) --------------------------------
\usepackage[hidelinks,breaklinks=true]{hyperref}
\hypersetup{
  pdftitle={EduCurator AI: desarrollo y evaluación en curso de un sistema agéntico para la curación asistida de contenidos educativos},
  pdfauthor={Correa Restrepo, Manco Henao, Pico Garzón, Méndez Roa},
  pdfsubject={Proyecto de investigación en curso},
  pdfkeywords={curación de contenidos educativos, RAG, agentes LLM, human-in-the-loop, design science research}
}
\usepackage[spanish]{cleveref}

% --- Comandos propios del proyecto -----------------------------------
% Marca de pendiente: visible en el PDF para el equipo, se desactiva
% cambiando \pendienteactivotrue por \pendienteactivofalse.
\newif\ifpendienteactivo
\pendienteactivotrue
\newcommand{\pend}[1]{\ifpendienteactivo\textcolor{red}{\small\textsf{[PENDIENTE: #1]}}\fi}

% Etiquetas del criterio de clasificación de párrafos (sección 24 del
% documento de ajuste de enfoque).
\newcommand{\etLit}{\textcolor{ecAzul}{\textsuperscript{(L)}}}   % literatura
\newcommand{\etMvp}{\textcolor{ecAzul}{\textsuperscript{(A)}}}   % artefacto
\newcommand{\etEval}{\textcolor{ecAzul}{\textsuperscript{(E)}}}  % a evaluar
\newcommand{\etAbierto}{\textcolor{ecAzul}{\textsuperscript{(?)}}} % desconocido
